\chapter*{前言（重构版）}
\addcontentsline{toc}{chapter}{前言（重构版）}

这本书是《社会动力学与公理化体系》的结构重构版。它的前身已经回答了一个问题：地理、文化、经济、政治与法律，能不能被放进同一条清晰的推导链里？答案是能——但那次回答只完成了一半。

前身的结构是\textbf{回溯}的：从结果出发，问「为什么会这样」。沿着地理$\rightarrow$结构$\rightarrow$经济$\rightarrow$文化$\rightarrow$政治$\rightarrow$法律一路回溯，它确实给出了一致的解释。但作者在自审报告里承认了三处结构性的病：数学附录里「空间正交化」到「因果条件概率」的衔接断裂（两者并不同构）；四个领域被写成了并列的四章，像四本教科书的合订本；以及最根本的一条——\textbf{它没有预测能力}。一本只解释不预测的社会理论，就像只有后视镜的汽车：看得清来路，看不清前路。

于是有了这一版。新结构是\textbf{前向}的：从动机出发，到目标，到未来，到各因素的贡献，再到具体的通道。关键转向只有一句话——原书问「为什么是这个结果」，本书问「这个结果里，各因素各占多少」。前者是历史学的问法，后者是\textbf{预测学}的问法：只有能定量地说出各占多少，才能回答「如果当初没有 X，今天会怎样」，而反事实推演（counterfactual）正是预测的最小单元。

怎么读这一版：全书五部，是一条前向的链——动机与目标（Part I）、演化状态空间（Part II）、条件概率与贡献分解（Part III）、四个通道（Part IV）、预测与循环（Part V）。每一章的开头都是上一章遗留的问题；每一个新概念都通过「加入一个约束」导出，不凭空引入。若暂时不想读公式，可以先跳过所有「数学建模」盒子——它们是落脚点，不是必经之路。若想知道一个结论有多可靠，留意红色的强调：那是对「这是猜想，不是定理」的诚实标注。

还要交代一句诚实的话：本书给出的是\textbf{结构}，不是\textbf{预言}。它能告诉你一个社会事件「依赖哪几个因素、各占多少」，不能告诉你「明天会发生什么」。方法会在哪里失败，第13章列了一张失败模式清单——那是本书可信度的一部分，不是它的污点。

原书的材料没有被丢弃，全部章节都在新结构里找到了新位置——这是重构，不是重写。它依然在生长，但已经足够诚恳。
